%% ============================================================
%%  Rapport — Projet Felidae Classifier
%%  ESGI — 3IABD2
%%  Compilez avec : pdflatex rapport_felidae.tex
%% ============================================================
\documentclass[12pt, a4paper]{report}

%% ---- Encodage & langue ------------------------------------
\usepackage[utf8]{inputenc}
\usepackage[T1]{fontenc}
\usepackage[french]{babel}

%% ---- Mise en page -----------------------------------------
\usepackage[top=2.5cm, bottom=2.5cm, left=2.5cm, right=2.5cm]{geometry}
\usepackage{setspace}
\onehalfspacing

%% ---- Mathématiques ----------------------------------------
\usepackage{amsmath}
\usepackage{amssymb}

%% ---- Code source ------------------------------------------
\usepackage{listings}
\usepackage{xcolor}

\definecolor{codegreen}{rgb}{0,0.6,0}
\definecolor{codegray}{rgb}{0.5,0.5,0.5}
\definecolor{codepurple}{rgb}{0.58,0,0.82}
\definecolor{backcolour}{rgb}{0.97,0.97,0.97}

\lstdefinelanguage{Rust}{
  keywords={fn, let, mut, pub, struct, impl, use, for, if, else, return,
            usize, f32, Vec, Self, pub, mod, crate, in, true, false},
  keywordstyle=\color{blue}\bfseries,
  ndkeywords={Matrix, LinearClassifier},
  ndkeywordstyle=\color{codepurple}\bfseries,
  comment=[l]{//},
  commentstyle=\color{codegreen}\itshape,
  stringstyle=\color{orange},
  morestring=[b]",
  sensitive=true
}

\lstset{
  language=Rust,
  backgroundcolor=\color{backcolour},
  basicstyle=\ttfamily\small,
  breaklines=true,
  captionpos=b,
  keepspaces=true,
  numbers=left,
  numberstyle=\tiny\color{codegray},
  numbersep=5pt,
  showspaces=false,
  showstringspaces=false,
  tabsize=4,
  frame=single,
  rulecolor=\color{codegray},
}

%% ---- Graphiques & tableaux --------------------------------
\usepackage{graphicx}
\usepackage{booktabs}
\usepackage{array}
\usepackage{float}

%% ---- Liens ------------------------------------------------
\usepackage[hidelinks]{hyperref}

%% ---- En-têtes & pieds de page ----------------------------
\usepackage{fancyhdr}
\pagestyle{fancy}
\fancyhf{}
\rhead{Felidae Classifier}
\lhead{ESGI — 3IABD2}
\rfoot{\thepage}

%% ---- Titre ------------------------------------------------
\title{
    \vspace{2cm}
    \Huge\textbf{Felidae Classifier} \\[0.5cm]
    \Large Implémentation d'algorithmes de classification \\
    en Rust sans bibliothèque de machine learning \\[1cm]
    \large ESGI — 3\textsuperscript{ème} année IABD \\[0.5cm]
    \normalsize Projet — Intelligence Artificielle
    \vspace{1cm}
}
\author{Léo}
\date{\today}

%% ============================================================
\begin{document}

\maketitle
\thispagestyle{empty}

\tableofcontents
\newpage

%% ============================================================
\chapter{Introduction}
%% ============================================================

\section{Contexte et objectifs}

% TODO — Décrire le contexte du projet : classification d'images de félins
% (chats, lions, guépards) en Rust sans bibliothèque de machine learning.

Ce projet a pour objectif de concevoir et d'implémenter from scratch,
en langage Rust, une bibliothèque de classification d'images de félins
(chats, lions et guépards). L'implémentation couvre plusieurs modèles
et algorithmes vus en cours, depuis le modèle linéaire jusqu'au SVM avec
noyau RBF, en passant par le PMC et le réseau RBF.

\section{Architecture du projet}

% TODO — Décrire brièvement la structure des crates et modules.

Le projet est structuré comme une bibliothèque Rust (\texttt{felidae\_classifier})
exposant plusieurs modules indépendants :

\begin{itemize}
    \item \texttt{tensor/} — opérations matricielles de base
    \item \texttt{models/} — modèles de classification (linéaire, PMC, RBFN, SVM)
    \item \texttt{training/} — boucle d'entraînement, optimiseurs, fonctions de perte
    \item \texttt{data/} — chargement et augmentation des données
    \item \texttt{features/} — normalisation et extraction de features
    \item \texttt{io/} — sauvegarde et chargement des modèles
    \item \texttt{server/} — système client/serveur pour l'inférence
\end{itemize}

\section{Technologies utilisées}

\begin{table}[H]
\centering
\begin{tabular}{ll}
\toprule
\textbf{Technologie} & \textbf{Rôle} \\
\midrule
Rust 2021 & Langage principal \\
\texttt{image 0.25} & Décodage et redimensionnement des images \\
\texttt{serde / serde\_json 1} & Sérialisation des modèles \\
\texttt{rand 0.10} & Initialisation des poids \\
\texttt{rayon 1} & Parallélisation du chargement des données \\
\texttt{tokio 1 / axum 0.8} & Serveur d'inférence asynchrone \\
\bottomrule
\end{tabular}
\caption{Dépendances principales du projet}
\end{table}

%% ============================================================
\chapter{Rendu 1 — Modèles linéaires et PMC}
%% ============================================================

%% ------------------------------------------------------------
\section{Module tensor}
%% ------------------------------------------------------------

\subsection{Rôle et conception}

% TODO — Expliquer pourquoi ce module est la fondation de tout le projet.

Le module \texttt{tensor/mod.rs} constitue la fondation de l'ensemble
du projet. Il définit la structure \texttt{Matrix}, un tableau 2D stocké
en mémoire de manière contiguë (\textit{row-major}), identique à la
convention des tableaux C.

\begin{lstlisting}[caption={Structure Matrix}]
pub struct Matrix {
    pub data: Vec<f32>,
    pub rows: usize,
    pub cols: usize,
}
\end{lstlisting}

L'accès à l'élément $(i, j)$ se calcule par :
\[
    \text{index} = i \times \text{cols} + j
\]

\subsection{Opérations implémentées}

% TODO — Lister et expliquer les opérations (dot, transpose, add, sub, scale, map, add_bias_row).

Les opérations suivantes ont été implémentées :

\begin{table}[H]
\centering
\begin{tabular}{lll}
\toprule
\textbf{Méthode} & \textbf{Équivalent numpy} & \textbf{Utilisation dans le projet} \\
\midrule
\texttt{dot} & \texttt{a @ b} & Passage avant de tous les modèles \\
\texttt{transpose} & \texttt{a.T} & Calcul des gradients (rétropropagation) \\
\texttt{add / sub} & \texttt{a + b} & Mise à jour des poids \\
\texttt{scale} & \texttt{a * k} & Application du taux d'apprentissage \\
\texttt{mul\_elementwise} & \texttt{a * b} & Rétropropagation dans le PMC \\
\texttt{map} & \texttt{np.vectorize(f)(a)} & Application des fonctions d'activation \\
\texttt{add\_bias\_row} & broadcast \texttt{+ b} & Ajout du biais dans le passage avant \\
\bottomrule
\end{tabular}
\caption{Opérations du module tensor}
\end{table}

\subsection{Tests unitaires}

% TODO — Décrire les tests et leurs résultats.

%% ------------------------------------------------------------
\section{Cas de tests synthétiques}
%% ------------------------------------------------------------

\subsection{Motivation}

Avant d'appliquer les modèles à des images réelles, des jeux de données
synthétiques ont été construits pour valider le comportement des modèles
dans des conditions contrôlées.

\subsection{Jeu de données 1 — Linéairement séparable}

% TODO — Décrire les 3 clusters (chat, lion, guépard) et inclure une visualisation.

Les trois classes sont représentées par des clusters de points 2D
bien séparés dans l'espace :

\begin{itemize}
    \item Classe 0 (chat) — coin inférieur gauche, centré sur $(0.15, 0.16)$
    \item Classe 1 (lion) — zone supérieure, centré sur $(0.52, 0.90)$
    \item Classe 2 (guépard) — coin inférieur droit, centré sur $(0.90, 0.14)$
\end{itemize}

% TODO — Ajouter une figure ici (scatter plot des points)
% \begin{figure}[H]
%     \centering
%     \includegraphics[width=0.6\textwidth]{figures/dataset1.png}
%     \caption{Jeu de données linéairement séparable}
% \end{figure}

\subsection{Jeu de données 2 — Non linéairement séparable (cas KO)}

% TODO — Décrire la disposition XOR-like et pourquoi le modèle linéaire échoue.

Le second jeu de données présente une disposition de type XOR : les
classes 0 et 1 sont disposées en diagonales opposées, rendant toute
séparation linéaire impossible.

%% ------------------------------------------------------------
\section{Modèle linéaire}
%% ------------------------------------------------------------

\subsection{Principe mathématique}

Le modèle linéaire calcule un score par classe selon :
\[
    \mathbf{O} = \mathbf{X} \cdot \mathbf{W} + \mathbf{b}
\]

où $\mathbf{X} \in \mathbb{R}^{N \times d}$, $\mathbf{W} \in \mathbb{R}^{d \times C}$,
$\mathbf{b} \in \mathbb{R}^{1 \times C}$, $N$ est le nombre d'échantillons,
$d$ le nombre de features et $C$ le nombre de classes.

\subsection{Fonction Softmax}

Les scores bruts (logits) sont convertis en probabilités via la fonction softmax :
\[
    \text{softmax}(\mathbf{z})_j = \frac{e^{z_j - \max(\mathbf{z})}}{\sum_{k} e^{z_k - \max(\mathbf{z})}}
\]

La soustraction du maximum est une astuce de stabilité numérique qui
évite le débordement de la fonction exponentielle.

\subsection{Fonction de perte — Entropie croisée}

\[
    \mathcal{L} = -\frac{1}{N} \sum_{i=1}^{N} \log\left(\hat{p}_{i, y_i}\right)
\]

où $\hat{p}_{i, y_i}$ est la probabilité prédite pour la vraie classe
de l'échantillon $i$.

\subsection{Rétropropagation}

Le gradient par rapport aux poids est calculé par :
\[
    \frac{\partial \mathcal{L}}{\partial \mathbf{W}} =
    \frac{1}{N} \mathbf{X}^{\top} \cdot (\hat{\mathbf{P}} - \mathbf{Y}_{\text{onehot}})
\]

\subsection{Résultats sur les données synthétiques}

% TODO — Remplir avec les résultats réels après exécution.

\begin{table}[H]
\centering
\begin{tabular}{lcc}
\toprule
\textbf{Jeu de données} & \textbf{Précision finale} & \textbf{Perte finale} \\
\midrule
Linéairement séparable & TODO \% & TODO \\
Non linéairement séparable & TODO \% & TODO \\
\bottomrule
\end{tabular}
\caption{Résultats du modèle linéaire sur les données synthétiques}
\end{table}

% TODO — Commenter les résultats observés.

%% ------------------------------------------------------------
\section{Transformation non linéaire pour les cas KO}
%% ------------------------------------------------------------

\subsection{Principe}

% TODO — Expliquer la transformation (x, y) → (x², y², x·y) et son effet.

Lorsque le modèle linéaire échoue sur un cas KO, une transformation
non linéaire est appliquée à l'espace d'entrée avant de réentraîner
le classifieur. Par exemple :
\[
    [x, y] \longrightarrow [x^2,\ y^2,\ x \cdot y]
\]

Cette transformation projette les données dans un espace de dimension
supérieure où elles deviennent linéairement séparables.

\subsection{Résultats après transformation}

% TODO — Comparer avant/après transformation.

\begin{table}[H]
\centering
\begin{tabular}{lcc}
\toprule
\textbf{Modèle} & \textbf{Précision (cas KO)} & \textbf{Commentaire} \\
\midrule
Linéaire brut & TODO \% & \\
Linéaire + transformation & TODO \% & \\
\bottomrule
\end{tabular}
\caption{Effet de la transformation non linéaire sur les cas KO}
\end{table}

%% ------------------------------------------------------------
\section{Perceptron Multi-Couches (PMC)}
%% ------------------------------------------------------------

\subsection{Architecture}

% TODO — Décrire l'architecture retenue (nombre de couches, neurones, activation).

\begin{table}[H]
\centering
\begin{tabular}{lll}
\toprule
\textbf{Couche} & \textbf{Taille} & \textbf{Activation} \\
\midrule
Entrée & TODO & — \\
Cachée 1 & TODO & ReLU \\
Cachée 2 & TODO & ReLU \\
Sortie & 3 & Softmax \\
\bottomrule
\end{tabular}
\caption{Architecture du PMC}
\end{table}

\subsection{Initialisation des poids}

% TODO — Décrire l'initialisation He et pourquoi elle est adaptée à ReLU.

\subsection{Rétropropagation}

% TODO — Décrire le calcul des gradients couche par couche.

\subsection{Résultats sur les données synthétiques}

% TODO — Remplir avec les résultats réels.

\begin{table}[H]
\centering
\begin{tabular}{lcc}
\toprule
\textbf{Jeu de données} & \textbf{Précision finale} & \textbf{Perte finale} \\
\midrule
Linéairement séparable & TODO \% & TODO \\
Non linéairement séparable & TODO \% & TODO \\
\bottomrule
\end{tabular}
\caption{Résultats du PMC sur les données synthétiques}
\end{table}

%% ------------------------------------------------------------
\section{Comparaison des modèles — Rendu 1}
%% ------------------------------------------------------------

% TODO — Comparer linéaire, linéaire+transformation et PMC sur les deux jeux de données.

\begin{table}[H]
\centering
\begin{tabular}{lccc}
\toprule
\textbf{Modèle} & \textbf{Dataset 1} & \textbf{Dataset 2} & \textbf{Dataset réel} \\
\midrule
Linéaire & TODO \% & TODO \% & TODO \% \\
Linéaire + transformation & TODO \% & TODO \% & TODO \% \\
PMC & TODO \% & TODO \% & TODO \% \\
\bottomrule
\end{tabular}
\caption{Comparaison des modèles — Rendu 1}
\end{table}

%% ------------------------------------------------------------
\section{Application sur le dataset réel}
%% ------------------------------------------------------------

\subsection{Description du dataset}

% TODO — Décrire le dataset (nombre d'images par classe, résolution, source).

\subsection{Pipeline de données}

% TODO — Décrire le chargement, redimensionnement (64×64), normalisation.

\subsection{Résultats}

% TODO — Remplir avec les résultats réels sur une portion du dataset.

%% ============================================================
\chapter{Rendu 2 — RBFN, SVM et système client/serveur}
%% ============================================================

%% ------------------------------------------------------------
\section{Réseau RBF (RBFN)}
%% ------------------------------------------------------------

\subsection{Principe}

% TODO — Décrire la couche cachée (centres), la fonction gaussienne et la couche de sortie.

La fonction de base radiale gaussienne est définie par :
\[
    \phi(\mathbf{x}, \mathbf{c}_k) = \exp\left(-\gamma \|\mathbf{x} - \mathbf{c}_k\|^2\right)
\]

\subsection{Sélection des centres}

% TODO — Décrire la méthode retenue (k-means ou aléatoire).

\subsection{Résultats}

% TODO — Remplir après implémentation.

%% ------------------------------------------------------------
\section{SVM avec noyau RBF}
%% ------------------------------------------------------------

\subsection{Principe}

% TODO — Décrire le SVM, la marge maximale et le noyau RBF.

Le noyau RBF est défini par :
\[
    K(\mathbf{x}, \mathbf{y}) = \exp\left(-\gamma \|\mathbf{x} - \mathbf{y}\|^2\right)
\]

\subsection{Algorithme SMO}

% TODO — Décrire l'algorithme Sequential Minimal Optimization.

\subsection{Classification multiclasse}

% TODO — Décrire la stratégie One-vs-One pour 3 classes.

\subsection{Résultats}

% TODO — Remplir après implémentation.

%% ------------------------------------------------------------
\section{Comparaison globale de tous les modèles}
%% ------------------------------------------------------------

% TODO — Tableau comparatif final sur le dataset complet.

\begin{table}[H]
\centering
\begin{tabular}{lcccc}
\toprule
\textbf{Modèle} & \textbf{Dataset synth. 1} & \textbf{Dataset synth. 2} & \textbf{Dataset réel} & \textbf{Temps} \\
\midrule
Linéaire & & & & \\
Linéaire + transf. & & & & \\
PMC & & & & \\
RBFN & & & & \\
SVM RBF & & & & \\
\bottomrule
\end{tabular}
\caption{Comparaison globale de tous les modèles}
\end{table}

%% ------------------------------------------------------------
\section{Système client/serveur}
%% ------------------------------------------------------------

\subsection{Architecture}

% TODO — Décrire l'architecture HTTP (axum + tokio), les endpoints, le protocole.

\subsection{Sauvegarde et chargement des modèles}

% TODO — Décrire le format de sérialisation (JSON) et les fonctions save/load.

\subsection{Démonstration}

% TODO — Décrire une démonstration end-to-end : entraînement → sauvegarde → inférence via API.

%% ============================================================
\chapter{Conclusion}
%% ============================================================

% TODO — Synthèse des résultats, difficultés rencontrées, perspectives.

\section{Bilan}

\section{Difficultés rencontrées}

\section{Perspectives}

%% ============================================================
\end{document}
